\begin{theorem}[Field affine inverse is a certified affine map]
\label{thm:field-affine-inverse-is-certified-affine}
For a BEDC history-level $\FieldUp+\CommRingUp$ candidate under the hypotheses
of \autoref{thm:field-affine-two-sided-automorphism-composition}, let
$p_a:\mathsf{NonZero}_F(a)$ and $p_b:\mathsf{NonZero}_F(b)$ be the two
nonzero witnesses for
$$
T(x)=\mathsf{add}(\mathsf{mul}(\mathsf{mul}(a,x),b),d).
$$
For each $c:\mathsf{BHist}$ put
$$
s(c)=
\mathsf{mul}(\mathsf{inv}(a,p_a),
\mathsf{mul}(\mathsf{add}(c,\mathsf{neg}(d)),\mathsf{inv}(b,p_b))).
$$
Let
$A:=\mathsf{inv}(a,p_a)$ and $B:=\mathsf{inv}(b,p_b)$, and suppose
$p_A:\mathsf{NonZero}_F(A)$ and $p_B:\mathsf{NonZero}_F(B)$ are explicit
nonzero witnesses for the inverse coefficients. Define
$$
D:=\mathsf{mul}(\mathsf{mul}(A,\mathsf{neg}(d)),B)
$$
and
$$
S(c):=\mathsf{add}(\mathsf{mul}(\mathsf{mul}(A,c),B),D).
$$
Then every $c:\mathsf{BHist}$ satisfies $s(c)\hsame S(c)$, and with witnesses
$p_A,p_B$ the map $S$ is a certified two-sided affine automorphism by
\autoref{thm:field-affine-two-sided-certified-automorphism}.
\end{theorem}
\begin{proof}
Expand $s(c)$ and apply the distributivity field from $\CommRingUp$ to the
inner sum $c+\mathsf{neg}(d)$ after left multiplication by $A$. This separates
the term containing $c$ from the translated term containing $\mathsf{neg}(d)$.

Next reassociate the right multiplication by $B$. Multiplication associativity
identifies the $c$-term with $\mathsf{mul}(\mathsf{mul}(A,c),B)$ and identifies
the translated term with
$$
\mathsf{mul}(\mathsf{mul}(A,\mathsf{neg}(d)),B).
$$
The definition of $D$ then gives $s(c)\hsame S(c)$ for every
$c:\mathsf{BHist}$.

It remains to certify the displayed normal form. Apply
\autoref{thm:field-affine-two-sided-certified-automorphism} with left factor
$A$, right factor $B$, offset $D$, and the explicit inverse-coefficient
witnesses $p_A,p_B$. The resulting inverse-pair and classifier-iff clauses are
exactly the certified affine automorphism structure for $S$.
\end{proof}

\begin{theorem}[Field affine inverse has derived inverse-coefficient witnesses]
\label{thm:field-affine-inverse-witnesses-derived}
For a BEDC history-level $\FieldUp+\CommRingUp$ candidate under the hypotheses
of \autoref{thm:field-affine-inverse-is-certified-affine} before the
inverse-coefficient witnesses are supplied, also assume zero-left
multiplication $\mathsf{mul}(0,x)\hsame 0$, one-apartness
$\omega:(1\hsame 0)\to\mathsf{False}$, and conversion maps
$\alpha_x:((x\hsame 0)\to\mathsf{False})\to\mathsf{NonZero}_F(x)$.
Let $p_a:\mathsf{NonZero}_F(a)$ and $p_b:\mathsf{NonZero}_F(b)$ be the
nonzero witnesses for
$$
T(x)=\mathsf{add}(\mathsf{mul}(\mathsf{mul}(a,x),b),d).
$$
For each $c:\mathsf{BHist}$ put
$$
s(c)=
\mathsf{mul}(\mathsf{inv}(a,p_a),
\mathsf{mul}(\mathsf{add}(c,\mathsf{neg}(d)),\mathsf{inv}(b,p_b))).
$$
Let $A:=\mathsf{inv}(a,p_a)$ and $B:=\mathsf{inv}(b,p_b)$. Then there are
derived witnesses $p_A:\mathsf{NonZero}_F(A)$ and
$p_B:\mathsf{NonZero}_F(B)$ such that, for
$$
D:=\mathsf{mul}(\mathsf{mul}(A,\mathsf{neg}(d)),B)
$$
and
$$
S(c):=\mathsf{add}(\mathsf{mul}(\mathsf{mul}(A,c),B),D),
$$
every $c:\mathsf{BHist}$ satisfies $s(c)\hsame S(c)$, and $S$ is a certified
two-sided affine automorphism.
\end{theorem}
\begin{proof}
Apply \autoref{thm:field-inverse-nonzero-from-one-apartness} to $p_a$ with
the displayed zero-left multiplication law, one-apartness witness $\omega$, and
conversion map $\alpha_A$. This produces
$p_A:\mathsf{NonZero}_F(\mathsf{inv}(a,p_a))$, which is
$\mathsf{NonZero}_F(A)$ by the definition of $A$.

Apply the same theorem to $p_b$ with the same structural data and conversion
map $\alpha_B$. This gives
$p_B:\mathsf{NonZero}_F(\mathsf{inv}(b,p_b))$, hence
$\mathsf{NonZero}_F(B)$.

With $p_A$ and $p_B$ now available, invoke
\autoref{thm:field-affine-inverse-is-certified-affine}. Its normal-form clause
gives $s(c)\hsame S(c)$ for every $c:\mathsf{BHist}$, and its certification
clause is exactly the certified two-sided affine automorphism structure for
$S$.
\end{proof}

Per-object certificate: traditional target $\FieldUp$ affine inverse witness
derivation; BEDC source $\FieldUp+\CommRingUp$ over $\mathsf{BHist}$; carrier
$\mathsf{BHist}$; classifier $\hsame$; stability obligations zero-left
multiplication, one-apartness, inverse laws, and affine normal-form
certification; ledger the original coefficient witnesses and the derived
inverse-coefficient witnesses; core theorem
\autoref{thm:field-affine-inverse-witnesses-derived}; exactness predicate
inverse affine normal form with derived witnesses; standard bridge ordinary
affine automorphism group; status $\mathsf{ProofTarget}$.

Per-object certificate: traditional target inverse closure for affine
automorphisms; BEDC source $\FieldUp+\CommRingUp$ over $\mathsf{BHist}$;
carrier $\mathsf{BHist}$; classifier $\hsame$; stability obligations affine
equation exactness, distributivity, associativity, inverse laws, and nonzero
witnesses for inverse coefficients; ledger the original and inverse coefficient
witnesses plus offset; core theorem
\autoref{thm:field-affine-inverse-is-certified-affine}; exactness predicate
inverse affine normal form; standard bridge ordinary affine automorphism group;
status $\mathsf{ProofTarget}$.

\begin{theorem}[Field affine explicit inverse is a certified automorphism]
\label{thm:field-affine-explicit-inverse-certified-automorphism}
Under the hypotheses and notation of
\autoref{thm:field-affine-two-sided-certified-automorphism}, define
$$
T(x)=\mathsf{add}(\mathsf{mul}(\mathsf{mul}(a,x),b),d)
$$
and, for each $r:\mathsf{BHist}$, define
$$
s(r)=
\mathsf{mul}(\mathsf{inv}(a,p_a),
\mathsf{mul}(\mathsf{add}(r,\mathsf{neg}(d)),\mathsf{inv}(b,p_b))).
$$
Then every $c,e,x:\mathsf{BHist}$ satisfies
$$
s(T(x))\hsame x,
\qquad
T(s(c))\hsame c,
\qquad
s(c)\hsame s(e)\Longleftrightarrow c\hsame e.
$$
\end{theorem}
\begin{proof}
Apply
\autoref{thm:field-affine-two-sided-certified-automorphism} with witnesses
$p_a,p_b$ and offset $d$. Its inverse-pair clauses give
$s(T(x))\hsame x$ and $T(s(c))\hsame c$ directly.

For the forward direction of the inverse-map classifier iff, assume
$c\hsame e$. Compose $T(s(c))\hsame c$, the assumed equality, and the symmetry
of $T(s(e))\hsame e$ to obtain
$$
T(s(c))\hsame T(s(e)).
$$
Reflect this equality through the certified classifier iff for $T$ to conclude
$s(c)\hsame s(e)$.

For the reverse direction, assume $s(c)\hsame s(e)$. The
classifier-preservation direction for $T$ gives
$$
T(s(c))\hsame T(s(e)).
$$
Compose the symmetry of $T(s(c))\hsame c$, the displayed equality, and
$T(s(e))\hsame e$ to obtain $c\hsame e$.
\end{proof}
\leanchecked{BEDC.Derived.FieldUp.field\_affine\_explicit\_inverse\_certified\_automorphism}

Per-object certificate: traditional target inverse affine automorphism; BEDC
source $\FieldUp$ over $\mathsf{BHist}$; carrier $\mathsf{BHist}$; classifier
$\hsame$; stability obligations certified affine automorphism,
inverse-pair clauses, and classifier iff; ledger the two nonzero witnesses and
offset; core theorem
\autoref{thm:field-affine-explicit-inverse-certified-automorphism}; exactness
predicate inverse-map classifier iff; standard bridge deferred; status
$\mathsf{ProofTarget}$.

\begin{theorem}[Field affine explicit inverse is classifier-congruent]
\label{thm:field-affine-explicit-inverse-classifier-congruence}
Under the hypotheses and notation of
\autoref{thm:field-affine-two-sided-certified-automorphism}, define
$$
T(x)=\mathsf{add}(\mathsf{mul}(\mathsf{mul}(a,x),b),d)
$$
and
$$
s(c)=
\mathsf{mul}(\mathsf{inv}(a,p_a),
\mathsf{mul}(\mathsf{add}(c,\mathsf{neg}(d)),\mathsf{inv}(b,p_b))).
$$
If $c\hsame c'$, then
$$
s(c)\hsame s(c').
$$
\end{theorem}
\begin{proof}
The left inverse clauses of
\autoref{thm:field-affine-two-sided-certified-automorphism} give
$T(s(c))\hsame c$ and $T(s(c'))\hsame c'$. Compose the first equality with
$c\hsame c'$ and with the symmetry of the second equality. This produces
$$
T(s(c))\hsame T(s(c')).
$$

The classifier iff in
\autoref{thm:field-affine-two-sided-certified-automorphism} reflects equality
of $T$-images back to equality of inputs. Applying that reflection direction to
the displayed classification yields $s(c)\hsame s(c')$.
\end{proof}

\begin{theorem}[Field affine pointwise classifier gives operation congruence]
\label{thm:field-affine-pointwise-classifier-operation-congruence}
For a BEDC history-level $\FieldUp+\CommRingUp$ candidate under the hypotheses
of \autoref{thm:field-affine-two-sided-certified-automorphism},
\autoref{thm:field-affine-two-sided-automorphism-composition}, and
\autoref{thm:field-affine-inverse-is-certified-affine}, let $T,T',U,U'$ be
certified displayed affine maps over $\mathsf{BHist}$. Write
$T^{-1}_{\mathrm{exp}}$ for the displayed explicit inverse formula supplied by
\autoref{thm:field-affine-two-sided-certified-automorphism}, and similarly for
the other certified maps. Suppose $\equiv_{\mathrm{aff}}$ is the pointwise
classifier of \autoref{def:field-affine-automorphism-package}. Then the
following clauses hold.
\begin{enumerate}
\item If $T\equiv_{\mathrm{aff}}T'$ and $U\equiv_{\mathrm{aff}}U'$, then
$$
U\circ T\equiv_{\mathrm{aff}}U'\circ T'.
$$
\item If $T\equiv_{\mathrm{aff}}U$, then
$$
T^{-1}_{\mathrm{exp}}\equiv_{\mathrm{aff}}U^{-1}_{\mathrm{exp}}.
$$
\end{enumerate}
\end{theorem}
\begin{proof}
For composition, fix $x:\mathsf{BHist}$. The assumption
$T\equiv_{\mathrm{aff}}T'$ gives $T(x)\hsame T'(x)$. The classifier-preservation
direction of the certified affine automorphism structure for $U$ transports
this to
$$
U(T(x))\hsame U(T'(x)).
$$

The assumption $U\equiv_{\mathrm{aff}}U'$ at the point $T'(x)$ gives
$$
U(T'(x))\hsame U'(T'(x)).
$$
Transitivity of $\hsame$ gives $U(T(x))\hsame U'(T'(x))$ for every $x$, which
is exactly $U\circ T\equiv_{\mathrm{aff}}U'\circ T'$.

For inverse congruence, fix $q:\mathsf{BHist}$. The explicit inverse clauses
give $T(T^{-1}_{\mathrm{exp}}(q))\hsame q$ and
$U(U^{-1}_{\mathrm{exp}}(q))\hsame q$. The pointwise comparison
$T\equiv_{\mathrm{aff}}U$ at $U^{-1}_{\mathrm{exp}}(q)$ gives
$$
T(U^{-1}_{\mathrm{exp}}(q))\hsame U(U^{-1}_{\mathrm{exp}}(q)).
$$

Thus $T(U^{-1}_{\mathrm{exp}}(q))\hsame q$. The certified fiber singleton
theorem \autoref{thm:field-affine-certified-fiber-singleton}, applied to the
two displayed preimages of $q$ under $T$, gives
$T^{-1}_{\mathrm{exp}}(q)\hsame U^{-1}_{\mathrm{exp}}(q)$. Since $q$ was
arbitrary, the explicit inverse formulas are related by
$\equiv_{\mathrm{aff}}$.
\end{proof}

Per-object certificate: traditional target affine inverse stability; BEDC source
$\FieldUp$ over $\mathsf{BHist}$; carrier $\mathsf{BHist}$; classifier
$\hsame$; stability obligations certified affine automorphism, left inverse,
and classifier reflection; ledger the two nonzero witnesses and offset; core
theorem
\autoref{thm:field-affine-explicit-inverse-classifier-congruence}; exactness
predicate inverse-map classifier congruence; standard bridge ordinary affine
inverse congruence; status $\mathsf{ProofTarget}$.

Per-object certificate: traditional target affine automorphism group law; BEDC
source $\FieldUp$ over $\mathsf{BHist}$; carrier $\mathsf{BHist}$; classifier
$\hsame$; stability obligations two certified affine automorphisms, explicit
inverse, and classifier iff; ledger four nonzero factor witnesses and two
offsets; core theorem
\autoref{thm:field-affine-composite-reverse-explicit-inverse}; exactness
predicate two-sided inverse plus composite classifier iff; standard bridge
deferred; status $\mathsf{ProofTarget}$.

\begin{theorem}[Field affine automorphism composition is associative]
\label{thm:field-affine-automorphism-composition-associative}
For a BEDC history-level $\FieldUp+\CommRingUp$ candidate under the hypotheses of
\autoref{thm:field-affine-two-sided-automorphism-composition}, let
$p_a,p_b,p_c,p_e,p_g,p_i$ be nonzero witnesses for the six displayed factors in
the certified affine maps
$$
T_1(x)=\mathsf{add}(\mathsf{mul}(\mathsf{mul}(a,x),b),d),
$$
$$
T_2(y)=\mathsf{add}(\mathsf{mul}(\mathsf{mul}(c,y),e),f),
$$
and
$$
T_3(z)=\mathsf{add}(\mathsf{mul}(\mathsf{mul}(g,z),i),j).
$$
Let $\mathsf{ApartZero}_F(h)$ abbreviate $h\hsame0\to\bot$. Assume the
product-apartness conversion maps required to certify the two binary composites
and the two outer composites:
$$
\alpha_{21}^L:\mathsf{ApartZero}_F(\mathsf{mul}(c,a))\to
\mathsf{NonZero}_F(\mathsf{mul}(c,a)),
\qquad
\alpha_{21}^R:\mathsf{ApartZero}_F(\mathsf{mul}(b,e))\to
\mathsf{NonZero}_F(\mathsf{mul}(b,e)),
$$
$$
\alpha_{32}^L:\mathsf{ApartZero}_F(\mathsf{mul}(g,c))\to
\mathsf{NonZero}_F(\mathsf{mul}(g,c)),
\qquad
\alpha_{32}^R:\mathsf{ApartZero}_F(\mathsf{mul}(e,i))\to
\mathsf{NonZero}_F(\mathsf{mul}(e,i)),
$$
$$
\alpha_{3(21)}^L:\mathsf{ApartZero}_F(\mathsf{mul}(g,\mathsf{mul}(c,a)))\to
\mathsf{NonZero}_F(\mathsf{mul}(g,\mathsf{mul}(c,a))),
\qquad
\alpha_{3(21)}^R:\mathsf{ApartZero}_F(\mathsf{mul}(\mathsf{mul}(b,e),i))\to
\mathsf{NonZero}_F(\mathsf{mul}(\mathsf{mul}(b,e),i)),
$$
and
$$
\alpha_{(32)1}^L:\mathsf{ApartZero}_F(\mathsf{mul}(\mathsf{mul}(g,c),a))\to
\mathsf{NonZero}_F(\mathsf{mul}(\mathsf{mul}(g,c),a)),
\qquad
\alpha_{(32)1}^R:\mathsf{ApartZero}_F(\mathsf{mul}(b,\mathsf{mul}(e,i)))\to
\mathsf{NonZero}_F(\mathsf{mul}(b,\mathsf{mul}(e,i))).
$$
Put
$$
A_{21}=\mathsf{mul}(c,a),\quad
B_{21}=\mathsf{mul}(b,e),\quad
D_{21}=\mathsf{add}(\mathsf{mul}(\mathsf{mul}(c,d),e),f),
$$
$$
A_{32}=\mathsf{mul}(g,c),\quad
B_{32}=\mathsf{mul}(e,i),\quad
D_{32}=\mathsf{add}(\mathsf{mul}(\mathsf{mul}(g,f),i),j).
$$
The bracketing $T_3\circ(T_2\circ T_1)$ supplies
$$
C_L(x)=
\mathsf{add}(\mathsf{mul}(\mathsf{mul}(\mathsf{mul}(g,A_{21}),x),
\mathsf{mul}(B_{21},i)),
\mathsf{add}(\mathsf{mul}(\mathsf{mul}(g,D_{21}),i),j)),
$$
and the bracketing $(T_3\circ T_2)\circ T_1$ supplies
$$
C_R(x)=
\mathsf{add}(\mathsf{mul}(\mathsf{mul}(\mathsf{mul}(A_{32},a),x),
\mathsf{mul}(b,B_{32})),
\mathsf{add}(\mathsf{mul}(\mathsf{mul}(A_{32},d),B_{32}),D_{32})).
$$
Then every $x,q:\mathsf{BHist}$ satisfies
$$
C_L(x)\hsame C_R(x).
$$
Moreover, if $s_1,s_2,s_3$ are the reverse explicit inverse formulas attached
to $T_1,T_2,T_3$, if $s_{21}$ is the reverse explicit inverse for the certified
composite $T_2\circ T_1$, and if $s_{32}$ is the reverse explicit inverse for
$T_3\circ T_2$, then the two bracketed reverse inverse formulas classify
together:
$$
s_{21}(s_3(q))\hsame s_1(s_{32}(q)).
$$
\end{theorem}
\begin{proof}
First certify the inner composites. Apply
\autoref{thm:field-affine-two-sided-automorphism-composition} to $T_2\circ T_1$
using $\alpha_{21}^L,\alpha_{21}^R$, and apply the same theorem to
$T_3\circ T_2$ using $\alpha_{32}^L,\alpha_{32}^R$. The collected data are
respectively $A_{21},B_{21},D_{21}$ and $A_{32},B_{32},D_{32}$.

Next certify the two outer composites. Apply
\autoref{thm:field-affine-two-sided-automorphism-composition} to
$T_3\circ(T_2\circ T_1)$ using $\alpha_{3(21)}^L,\alpha_{3(21)}^R$, and apply it
to $(T_3\circ T_2)\circ T_1$ using $\alpha_{(32)1}^L,\alpha_{(32)1}^R$.
Both certified maps classify the same term $T_3(T_2(T_1(x)))$.

The collected left factors are
$\mathsf{mul}(g,\mathsf{mul}(c,a))$ and
$\mathsf{mul}(\mathsf{mul}(g,c),a)$, the collected right factors are
$\mathsf{mul}(\mathsf{mul}(b,e),i)$ and
$\mathsf{mul}(b,\mathsf{mul}(e,i))$. Multiplication associativity identifies
each pair.

It remains to compare the offsets. Expand $D_{21}$ and $D_{32}$. Distributivity
moves the middle offset $f$ through the final affine map, and multiplication
associativity collects the factors multiplying $d$. The two displayed offsets
therefore classify together, so the normal forms $C_L(x)$ and $C_R(x)$ classify
together for every $x$.

For the inverse formulas, apply
\autoref{thm:field-affine-composite-inverse-reverse-coherence} to
$T_2\circ T_1$ and then compose the resulting equality with $s_3(q)$. Apply the
same coherence theorem to $T_3\circ T_2$ at $q$, then transport through $s_1$.
Transitivity gives
$s_{21}(s_3(q))\hsame s_1(s_2(s_3(q)))\hsame s_1(s_{32}(q))$, which is the
displayed equality.
\end{proof}

Per-object certificate: traditional target affine automorphism associativity;
BEDC source $\FieldUp+\CommRingUp$ over $\mathsf{BHist}$; carrier
$\mathsf{BHist}$; classifier $\hsame$; stability obligations binary affine
composition, product apartness, distributivity, inverse coherence; ledger the
six factor witnesses, product witnesses, and three offsets; core theorem
\autoref{thm:field-affine-automorphism-composition-associative}; exactness
predicate associativity of certified affine composition; standard bridge
ordinary affine group law; status $\mathsf{ProofTarget}$.

\begin{theorem}[Field affine automorphism identity and inverse package]
\label{thm:field-affine-automorphism-identity-and-inverse-package}
For a BEDC history-level $\FieldUp$ candidate over $\mathsf{BHist}$ under the
hypotheses used by
\autoref{thm:field-affine-two-sided-certified-automorphism} and
\autoref{thm:field-affine-two-sided-automorphism-composition}, bind the unit
witness $p_1:\mathsf{NonZero}_F(1)$ and define
$$
I(x)=\mathsf{add}(\mathsf{mul}(\mathsf{mul}(1,x),1),0).
$$
For every certified two-sided affine map
$$
T(x)=\mathsf{add}(\mathsf{mul}(\mathsf{mul}(a,x),b),d)
$$
with witnesses $p_a:\mathsf{NonZero}_F(a)$ and
$p_b:\mathsf{NonZero}_F(b)$, every $x:\mathsf{BHist}$ satisfies
$$
I(x)\hsame x,
\qquad
I(T(x))\hsame T(x),
\qquad
T(I(x))\hsame T(x).
$$
For every second certified two-sided affine map
$$
U(z)=\mathsf{add}(\mathsf{mul}(\mathsf{mul}(c,z),e),f)
$$
with witnesses $p_c:\mathsf{NonZero}_F(c)$ and
$p_e:\mathsf{NonZero}_F(e)$, let $\mathsf{ApartZero}_F(h)$ abbreviate
$h\hsame0\to\bot$. Assume conversion maps
$$
\alpha_{ca}:\mathsf{ApartZero}_F(\mathsf{mul}(c,a))\to
\mathsf{NonZero}_F(\mathsf{mul}(c,a)),
\qquad
\alpha_{be}:\mathsf{ApartZero}_F(\mathsf{mul}(b,e))\to
\mathsf{NonZero}_F(\mathsf{mul}(b,e)).
$$
Let $q_{ca}$ and $q_{be}$ be the product-apartness witnesses supplied by
\autoref{thm:field-nonzero-factors-exclude-zero-product} for the pairs
$(p_c,p_a)$ and $(p_b,p_e)$, and set
$p_{ca}:=\alpha_{ca}(q_{ca})$ and $p_{be}:=\alpha_{be}(q_{be})$. Then the
composite $U(T(x))$ is certified by the collected coefficients
$$
A=\mathsf{mul}(c,a),
\qquad
B=\mathsf{mul}(b,e),
\qquad
D=\mathsf{add}(\mathsf{mul}(\mathsf{mul}(c,d),e),f)
$$
together with the product witnesses $p_{ca}$ and $p_{be}$.
Moreover, if $s_T$ and $s_U$ are the displayed inverse maps for $T$ and $U$,
then $s_T(s_U(q))$ is a two-sided inverse of $U(T(x))$ and agrees with the
collected-coefficient inverse of the composite indexed by $p_{ca}$ and
$p_{be}$. Hence the identity affine map, composition, inverse closure, and
classifier iff clauses form the affine automorphism package over
$\mathsf{BHist}$.
\end{theorem}
\begin{proof}
The identity clauses use the same $\FieldUp$ operations as the certified affine
maps. The multiplicative unit laws reduce $\mathsf{mul}(\mathsf{mul}(1,x),1)$
to $x$, and the additive zero law reduces the outer addition to $x$; substituting
$T(x)$ for $x$ gives $I(T(x))\hsame T(x)$. Classifier preservation for $T$ from
\autoref{thm:field-affine-two-sided-certified-automorphism} transports
$I(x)\hsame x$ through $T$, giving $T(I(x))\hsame T(x)$.

For composition, apply
\autoref{thm:field-affine-two-sided-automorphism-composition} to the two
certified affine maps using the displayed $\alpha_{ca}$ and $\alpha_{be}$ data.
The product-apartness witnesses $q_{ca}$ and $q_{be}$ supplied by
\autoref{thm:field-nonzero-factors-exclude-zero-product} are converted into
$p_{ca}$ and $p_{be}$, so the composition theorem supplies the collected
coefficients $A,B,D$, the certified composite indexed by those product
witnesses, and the classifier iff for $U\circ T$. The certified-automorphism
clauses are then read through
\autoref{thm:field-affine-two-sided-certified-automorphism}.

For inverse closure, use
\autoref{thm:field-affine-composite-reverse-explicit-inverse}. It states that
the reverse composite $s_T(s_U(q))$ is both a left preimage and a right inverse
for $U(T(x))$, while preserving the classifier iff. Finally,
\autoref{thm:field-affine-composite-inverse-reverse-coherence} identifies this
reverse composite with the inverse obtained directly from the collected
coefficients, so the inverse operation is closed under the same ledger data.
\end{proof}

Per-object certificate: traditional target affine automorphism group interface;
BEDC source $\FieldUp+\CommRingUp$ over $\mathsf{BHist}$; carrier
$\mathsf{BHist}$; classifier $\hsame$; stability obligations identity,
composition, inverse, and classifier iff; ledger factor witnesses, product
witnesses, inverse witnesses, and offsets; core theorem
\autoref{thm:field-affine-automorphism-identity-and-inverse-package};
exactness predicate identity plus inverse-closed composition; standard bridge
deferred; status $\mathsf{ProofTarget}$.

\begin{theorem}[Field affine automorphisms form a GroupUp-compatible package]
\label{thm:field-affine-automorphism-group-laws}
For a BEDC history-level $\FieldUp+\CommRingUp$ candidate over
$\mathsf{BHist}$ under the hypotheses used by
\autoref{thm:field-affine-two-sided-certified-automorphism} and
\autoref{thm:field-affine-two-sided-automorphism-composition}, let
$\mathsf{AffAut}_F$ be the carrier of certified two-sided affine maps. A member
is a tuple $(a,b,d,p_a,p_b)$ with $a,b,d:\mathsf{BHist}$,
$p_a:\mathsf{NonZero}_F(a)$, and $p_b:\mathsf{NonZero}_F(b)$, displaying
$$
T_{a,b,d}(x)=\mathsf{add}(\mathsf{mul}(\mathsf{mul}(a,x),b),d).
$$
The affine classifier relates two members when their displayed maps are
pointwise $\hsame$, exactly as in
\autoref{def:field-affine-automorphism-package}. The identity member has
coefficients $(1,1,0)$ with the unit nonzero witnesses. The product of
$U=(c,e,f,p_c,p_e)$ after
$T=(a,b,d,p_a,p_b)$ has collected coefficients
$$
A=\mathsf{mul}(c,a),
\qquad
B=\mathsf{mul}(b,e),
\qquad
D=\mathsf{add}(\mathsf{mul}(\mathsf{mul}(c,d),e),f).
$$
The inverse operation is represented by the explicit inverse formula certified
in \autoref{thm:field-affine-inverse-is-certified-affine}. With these data,
$\mathsf{AffAut}_F$ satisfies the $\GroupUp$ source fields: carrier
inhabitance, closure under composition and inverse, associativity under the
affine classifier, left and right identity laws, left and right inverse laws,
and classifier congruence for both composition and inverse.
\end{theorem}
\begin{proof}
Carrier inhabitance is supplied by the identity coefficients. The unit nonzero
witnesses and the identity clauses of
\autoref{thm:field-affine-automorphism-identity-and-inverse-package} show that
this member acts as the identity map on every $x:\mathsf{BHist}$, and hence is
the distinguished $\mathsf{AffAut}_F$ element.

Composition closure is the collected-coefficient construction of
\autoref{thm:field-affine-two-sided-automorphism-composition}. It supplies the
product nonzero witnesses for $A$ and $B$, the displayed affine normal form for
$U(T(x))$, and the classifier iff clauses for the composite. The certified
two-sided behavior of each displayed map is read through
\autoref{thm:field-affine-two-sided-certified-automorphism}.

Associativity is exactly the pointwise comparison from
\autoref{thm:field-affine-automorphism-composition-associative}. For three
certified affine members, the two parenthesized compositions collect to
classifier-equal displayed maps, so the pointwise affine classifier identifies
the two products.

The inverse operation uses
\autoref{thm:field-affine-automorphism-identity-and-inverse-package} for the
single-map identity and inverse clauses, together with
\autoref{thm:field-affine-inverse-is-certified-affine} for inverse closure.
These clauses give left and right inverse laws under the affine classifier.

Classifier congruence for composition and inverse is exactly
\autoref{thm:field-affine-pointwise-classifier-operation-congruence}. Hence the
pointwise classifier supplies the operation fields required by the
$\GroupUp$ package.
\end{proof}

Per-object certificate: traditional target affine automorphism $\GroupUp$
interface; BEDC source $\FieldUp+\CommRingUp$ over $\mathsf{BHist}$; carrier
$\mathsf{AffAut}_F$; classifier pointwise $\hsame$ on displayed affine maps;
stability obligations inhabitance, composition, inverse, associativity,
identity, inverse laws, and classifier congruence; ledger identity witnesses,
side-factor witnesses, product witnesses, reverse inverse formulas, and
offsets; core theorem
\autoref{thm:field-affine-automorphism-group-laws}; exactness predicate
$\GroupUp$ source-field satisfaction for certified affine automorphisms;
standard bridge ordinary affine automorphism group law; status
$\mathsf{ProofTarget}$.

\begin{theorem}[Field affine product is independent of product witnesses]
\label{thm:field-affine-product-witness-independence}
For a BEDC history-level $\FieldUp+\CommRingUp$ candidate over
$\mathsf{BHist}$ under the hypotheses of
\autoref{thm:field-affine-two-sided-automorphism-composition}, let
$T=(a,b,d,p_a,p_b)$ and $U=(c,e,f,p_c,p_e)$ be certified affine members of
$\mathsf{AffAut}_F$. Set
$$
A=\mathsf{mul}(c,a),
\qquad
B=\mathsf{mul}(b,e),
\qquad
D=\mathsf{add}(\mathsf{mul}(\mathsf{mul}(c,d),e),f).
$$
Let two choices of apartness-to-nonzero conversion data produce witnesses
$$
p_{ca},p'_{ca}:\mathsf{NonZero}_F(A),
\qquad
p_{be},p'_{be}:\mathsf{NonZero}_F(B).
$$
Form the two composite affine members with the same collected coefficients
$(A,B,D)$ and with witness pairs $(p_{ca},p_{be})$ and
$(p'_{ca},p'_{be})$. Then the affine classifier on $\mathsf{AffAut}_F$
relates these two composites: their displayed maps are pointwise $\hsame$,
and their reverse explicit inverse formulas are pointwise $\hsame$.
\end{theorem}
\leanchecked{BEDC.Derived.FieldUp.field\_affine\_product\_witness\_independence}
\begin{proof}
The displayed composite maps have identical coefficients $A$ and $B$ and the
same offset $D$. For each $x:\mathsf{BHist}$, both displays reduce to
$\mathsf{add}(\mathsf{mul}(\mathsf{mul}(A,x),B),D)$, so pointwise
classification is reflexive.

It remains to compare the proof-indexed inverse factors occurring in the
reverse explicit inverse formulas. Apply
\autoref{thm:field-inverse-congruence-from-apartness} to the reflexive
coefficient comparison $A\hsame A$ with witnesses $p_{ca}$ and $p'_{ca}$,
giving $\mathsf{inv}(A,p_{ca})\hsame\mathsf{inv}(A,p'_{ca})$. Apply the same
theorem to $B\hsame B$ with witnesses $p_{be}$ and $p'_{be}$, giving
$\mathsf{inv}(B,p_{be})\hsame\mathsf{inv}(B,p'_{be})$.

Multiplication congruence transports the first inverse comparison through the
left affine factor, multiplication congruence transports the second comparison
through the right affine factor, and addition congruence transports the shared
offset term. Therefore the two reverse explicit inverse formulas are
pointwise $\hsame$, which is exactly the second field of the affine
classifier.
\end{proof}

Per-object certificate: traditional target affine composition
well-definedness; BEDC source $\FieldUp+\CommRingUp$ over $\mathsf{BHist}$;
carrier $\mathsf{AffAut}_F$; classifier the affine classifier; stability
obligation proof-index independence for product witnesses; ledger the two
witness pairs $(p_{ca},p_{be})$ and $(p'_{ca},p'_{be})$; core theorem
\autoref{thm:field-affine-product-witness-independence}; exactness predicate
product operation descends to the affine classifier; standard bridge ordinary
affine automorphism composition; status $\mathsf{ProofTarget}$.

\paragraph{Capstone.}

\begin{theorem}[Field affine automorphism group certificate]
\label{thm:field-affine-automorphism-group-certificate}
For a BEDC history-level $\FieldUp+\CommRingUp$ candidate over
$\mathsf{BHist}$, let $\mathsf{AffAut}_F(T)$ mean that $T$ carries displayed
parameters $a,b,d:\mathsf{BHist}$, witnesses
$p_a:\mathsf{NonZero}_F(a)$ and $p_b:\mathsf{NonZero}_F(b)$, and the certified
two-sided affine automorphism clauses for
$$
T(x)=\mathsf{add}(\mathsf{mul}(\mathsf{mul}(a,x),b),d).
$$
Equip the affine classifier with pointwise $\hsame$ on all carried histories.
Then the identity affine map, the collected-coefficient composition law, and
the explicit inverse operation satisfy the $\GroupUp$ stability fields: left
and right identity, associativity of composition, inverse closure, left and
right inverse laws, and classifier congruence. Moreover, for certified members
$T=(a,b,d,p_a,p_b)$ and $U=(c,e,f,p_c,p_e)$ with fixed collected coefficients
$$
A=\mathsf{mul}(c,a),
\qquad
B=\mathsf{mul}(b,e),
\qquad
D=\mathsf{add}(\mathsf{mul}(\mathsf{mul}(c,d),e),f),
$$
any two product-witness choices $p_{ca},p_{be}$ and $p'_{ca},p'_{be}$ produce
composite affine members related by the affine pointwise classifier, including
their reverse explicit inverse formulas. Equivalently, the displayed affine
automorphism package determines the concrete $\GroupUp$ certificate for
$\mathsf{AffAut}_F$ over $\mathsf{BHist}$.
\end{theorem}
\begin{proof}
The identity and inverse data are exactly the package isolated in
\autoref{thm:field-affine-automorphism-identity-and-inverse-package}. The
identity affine map supplies both unit laws pointwise, and
\autoref{thm:field-affine-two-sided-explicit-inverse} supplies the displayed
inverse operation with its left and right inverse clauses.

Composition uses the collected coefficients from the same package, with the
two-sided certification read through
\autoref{thm:field-affine-two-sided-certified-automorphism}.

For proof-relevant product witnesses, fix collected coefficients $A$, $B$, and
$D$ and compare the two composites built from witness pairs
$(p_{ca},p_{be})$ and $(p'_{ca},p'_{be})$. Apply
\autoref{thm:field-affine-product-witness-independence}: it gives pointwise
$\hsame$ for the displayed composite maps and for their reverse explicit
inverse formulas, which discharges the witness-choice descent field of the
certificate.

The classifier is pointwise $\hsame$, so classifier congruence for composition
and inverse is obtained from
\autoref{thm:field-affine-pointwise-classifier-operation-congruence}.

Associativity is the normal-form statement of
\autoref{thm:field-affine-automorphism-composition-associative}. It identifies
the two parenthesized collected composites by the BEDC ring and field laws, and
therefore gives the associativity field of the $\GroupUp$ certificate.

The theorem
$\mathsf{BEDC.Derived.GroupUp.group\_stability\_certificate\_fields}$ is the
packaging result consumed by Lean: once the identity, composition,
associativity, inverse closure, inverse laws, and classifier congruence fields
are supplied, it repacks them as the public $\GroupUp$ stability certificate.
The preceding paragraphs provide precisely those fields for
$\mathsf{AffAut}_F$.
\end{proof}
